\documentclass[12pt]{exam}

\usepackage[T1]{fontenc}
\usepackage[utf8]{inputenc}
\usepackage[lithuanian]{babel}
\usepackage{lmodern}
\usepackage{graphicx}
\usepackage{xcolor}
\usepackage{tikz}
\usepackage[
    a4paper,
    left=1cm,
    right=1cm,
    top=2.5cm,
    bottom=1cm
]{geometry}
\usepackage{amsmath,amssymb}
\usepackage{tasks}

% Švelnus pilkas rėmelis mokinio įrašui
\newcommand{\irasolaukas}[1]{%
    \begingroup
    \setlength{\fboxrule}{0.3pt}%
    \setlength{\fboxsep}{3pt}%
    \fcolorbox{black!35}{white}{%
        \rule{0pt}{0.45cm}\hspace{#1}%
    }%
    \endgroup
}

% Sprendimo laukas:
% argumentas nurodo linijų skaičių;
% tarpas tarp linijų – 0.7 cm.
\newcommand{\sprendimas}[1]{%
    \par\smallskip
    \noindent\textit{Sprendimas:}
    \par\nobreak
    \vspace{0.1cm}
    \noindent
    \begin{tikzpicture}
        \path[use as bounding box]
            (0,0) rectangle (\linewidth,{-0.7*#1});
        \foreach \i in {1,...,#1} {
            \draw[
                draw=black!25,
                line width=0.3pt
            ]
                (0,{-0.7*\i})
                -- (\linewidth,{-0.7*\i});
        }
    \end{tikzpicture}%
    \par\smallskip
}

% Atsakymo laukelis
\newcommand{\atsakymas}[1][9]{%
    \par\smallskip
    \noindent\textit{Atsakymas:}
    \framebox[#1cm]{\rule{0pt}{0.5cm}}
    \par
}

% Antraštė
\pagestyle{headandfoot}
\runningheadrule

\firstpageheader
{%
    \small
    \begin{tabular}{@{}l@{\hspace{0.2cm}}l@{}}
        Mokinio kodas: & \irasolaukas{2.3cm} \\[0.15cm]
        Klasė:        & \irasolaukas{2.3cm}
    \end{tabular}%
}
{Kontrolinis darbas:\\Kvadratinės lygtys}
{Poseidono variantas}

\runningheader{Kvadratinės lygtys}{}{}
\firstpagefooter{}{}{}
\runningfooter{}{}{}

% Taškų žymėjimas
\pointsdroppedatright
\bracketedpoints
\setlength{\rightpointsmargin}{1.6cm}
\renewcommand{\points}{taškai}
\renewcommand{\pointname}{Taškai:}
\marginpointname{\ taškai}

\begin{document}

\begin{questions}

\question[3]
Išspręskite lygtį $7x=4+3x^2$.\droppoints
\sprendimas{4}
\atsakymas[3]

\question[2]
Sudarykite redukuotą kvadratinę lygtį, kurios sprendiniai būtų skaičiai $\frac{1}{3}$ ir $-5$.\droppoints
\sprendimas{2}
\atsakymas[3]

\question[7]
Išskaidykite dauginamaisiais:\droppoints
\begin{tasks}[label-width=1.2em](3)
\task $-2x^{2}+3x+5$
\task $x^{3}-5x^{2}-4x+20$
\task $x^{3}+2x^{2}-3$
\end{tasks}

\sprendimas{5}
\atsakymas[9]

\question[4]
Nustatykite trinario $P(x)=x^{2}+bx+9$ koeficiento $b$ reikšmę, jei jo šaknų skirtumas $8$.
\droppoints
\sprendimas{7}
\atsakymas[2]

\question[5]
Išspręskite lygtį $(x^{2}-5x)^{2}=2(x-2)(x-3)+12$.
\droppoints
\sprendimas{8}
\atsakymas[3]

\question[7]
Duotas kvadratinis trinaris $P(x)=-x^{2}+6x-3$. Žinoma, kad jo šaknys yra skaičiai $m$ ir $n$. Apskaičiuokite: \droppoints
\begin{tasks}(3)
    \task $m^{2}+n^{2}$
    \task $\frac{1}{m^{3}}+\frac{1}{n^{3}}$
    \task $m^{2}-n^{2}$, jei $m<n$
\end{tasks}
\sprendimas{6}
\atsakymas

\question[6]
Išspręskite lygtį $\frac{6-x}{x^{2}-4x+6}-\frac{4x^{2}-16x+24}{2x-12}=3$.
\droppoints
\sprendimas{9}
\atsakymas[3]

\question[6]
Kvadratinė lygtis $x^2-mx+9-m=0$ turi du realiuosius sprendinius $x_1,x_2$, tenkinančius sąlygą $x_1^2+x_2^2+x_1^2x_2^2=39$. Raskite $m$, neskaičiuodami pačių sprendinių. \droppoints
\sprendimas{13}
\atsakymas[2]
\end{questions}
\begin{center}
    \Large\textbf{Juodraštis}
\end{center}
\newpage
\begin{center}
    \Large\textbf{Juodraštis}
\end{center}

\end{document}
